\chapter{1958年黑吉辽卷}

\section{解答题}

\begin{problem}
已知：圆 \(O\) 内直交的二弦 \(AB\)，\(CD\) 相交于 \(P\). 求证：\(PA^2+PB^2+PC^2+PD^2\) 为一定值.

\bitmapfigure[max width=0.33\linewidth]{img/1958/hei_ji_liao/q1_fig1.png}
\end{problem}

\begin{solution}
设 \(PA=u\)，\(PB=v\)，\(PC=w\)，\(PD=z\)，其中 \(u\)，\(v\)，\(w\)，\(z>0\). 以 \(P\) 为原点，以 \(AB\)，\(CD\) 所在直线为坐标轴，则
\(A=\left(-u,0\right)\)，\(B=\left(v,0\right)\)，\(C=\left(0,w\right)\)，\(D=\left(0,-z\right)\). 设圆的方程为
\[
x^2+y^2-2\lambda x-2\mu y+c=0
\]
将 \(A\)，\(B\) 的坐标代入，得到关于 \(x\) 的二次方程
\(x^2-2\lambda x+c=0\)，其两个根为 \(-u\)，\(v\)，所以
\(2\lambda=v-u\)，\(c=-uv\). 再将 \(C\)，\(D\) 的坐标代入，得到关于 \(y\) 的二次方程
\(y^2-2\mu y+c=0\)，其两个根为 \(w\)，\(-z\)，所以 \(2\mu=w-z\)，\(c=-wz\). 因此 \(uv=wz\).

圆心为 \(\left(\lambda,\mu\right)\)，设半径为 \(R\)，则 \(R^2=\lambda^2+\mu^2-c\). 从而
\[
4R^2=\left(v-u\right)^2+\left(w-z\right)^2+4uv=u^2+v^2+w^2+z^2+2uv-2wz=u^2+v^2+w^2+z^2
\]
代回 \(u\)，\(v\)，\(w\)，\(z\) 的定义，得
\(PA^2+PB^2+PC^2+PD^2=4R^2\). 因而所求一定值为圆的直径的平方.
\end{solution}

\begin{problem}
工地上有沙堆，其下底面为矩形，长，宽各为 \(a\text{ 米}\)，\(b\text{ 米}\)，其各侧面都与地平面夹成 \(45^\circ\) 角，上，下底距离为 \(h\text{ 米}\)，求沙堆体积.

\bitmapfigure[max width=0.34\linewidth]{img/1958/hei_ji_liao/q2_fig1.png}
\end{problem}

\begin{solution}
在垂直于下底面一条边的截面内，侧面与地平面的夹角为 \(45^\circ\)，所以沙堆上升 \(h\text{ 米}\) 时，水平方向向内收进 \(h\text{ 米}\). 四个侧面都如此收进，故上底面的长，宽分别为
\(c=a-2h\)，\(d=b-2h\)，并且应有 \(a\ge2h\)，\(b\ge2h\).

将沙堆分解为一个底为 \(c\times d\)、高为 \(h\) 的直棱柱，沿 \(c\) 方向的两块底为直角三角形且两直角边均为 \(h\) 的三棱柱，沿 \(d\) 方向的另外两块同类三棱柱，以及四个底面为边长 \(h\) 的正方形、高为 \(h\) 的棱锥. 因此
\[
V=cdh+2\left(\frac12h^2c\right)+2\left(\frac12h^2d\right)+4\left(\frac13h^3\right)=h\left(a-2h\right)\left(b-2h\right)+h^2\left(a+b-4h\right)+\frac43h^3
\]
\[
V=abh-(a+b)h^2+\frac43h^3
\]
所以沙堆体积为 \(abh-(a+b)h^2+\frac43h^3\text{ 立方米}\).
\end{solution}

\begin{problem}
求根式：\(\sqrt[3]{a}\) 与 \(\sqrt{5-a}\) 的乘积.
\end{problem}

\begin{solution}
实数根式有意义的条件是 \(5-a\ge0\)，即 \(a\le5\). 当 \(0\le a\le5\) 时，\(\sqrt[3]{a}=\sqrt[6]{a^2}\)，于是
\[
\sqrt[3]{a}\sqrt{5-a}=\sqrt[6]{a^2}\sqrt[6]{\left(5-a\right)^3}=\sqrt[6]{a^2\left(5-a\right)^3}
\]
当 \(a<0\) 时，\(\sqrt[3]{a}=-\sqrt[6]{a^2}\)，而 \(\sqrt{5-a}=\sqrt[6]{(5-a)^3}>0\)，所以
\[
\sqrt[3]{a}\sqrt{5-a}=-\sqrt[6]{a^2\left(5-a\right)^3}
\]
因此，在实数范围内
\[
\sqrt[3]{a}\sqrt{5-a}=\begin{cases}
\sqrt[6]{a^2\left(5-a\right)^3},&0\le a\le5,\\
-\sqrt[6]{a^2\left(5-a\right)^3},&a<0.
\end{cases}
\]
\end{solution}

\begin{problem}
解方程：\(x^{1+\lg x}=20^{\log_{20}100}\).
\end{problem}

\begin{solution}
由对数的定义，\(20^{\log_{20}100}=100\)，所以原方程为 \(x^{1+\lg x}=100\). 令 \(t=\lg x\)，则定义域要求 \(x>0\)，且 \(x=10^t\). 于是
\[
10^{t\left(1+t\right)}=10^2
\]
从而 \(t(1+t)=2\)，即
\(t^2+t-2=0\)，\(\left(t+2\right)\left(t-1\right)=0\).
所以 \(t=1\) 或 \(t=-2\)，相应地 \(x=10\) 或 \(x=10^{-2}=\frac1{100}\). 两个数都大于零，并且分别代入可得 \(10^{1+1}=100\)，\(\left(10^{-2}\right)^{1-2}=10^2=100\)，故原方程的解为 \(x=10\)，\(x=\frac1{100}\).
\end{solution}

\begin{problem}
\begin{enumerate}
\item 解方程组：\(\begin{cases}
2x-5z=10-6a,\\
x-2y+3z=2a+1,\\
5x-6y+9z=23-10a.
\end{cases}\)
\item 若 \(x\) 为某直角三角形的斜边，\(y\) 及 \(z\) 为直角边，试确定 \(a\) 的值.
\end{enumerate}
\end{problem}

\begin{solution}
\begin{enumerate}
\item 将第二式乘以 \(3\)，得 \(3x-6y+9z=6a+3\). 用第三式减去这个等式，得到
\(2x=20-16a\)，所以 \(x=10-8a\). 代入第一式，得
\(2\left(10-8a\right)-5z=10-6a\)，所以 \(z=2-2a\). 再将 \(x\)，\(z\) 代入第二式，得
\[
10-8a-2y+3\left(2-2a\right)=2a+1
\]
即 \(y=\frac{15-16a}{2}\). 因而方程组的解为
\(x=10-8a\)，\(y=\frac{15-16a}{2}\)，\(z=2-2a\).

\item 三条边长必须为正，且由勾股定理
\[
(10-8a)^2=\left(\frac{15-16a}{2}\right)^2+(2-2a)^2
\]
两边乘以 \(4\)，并展开，得到
\[
400-640a+256a^2=225-480a+256a^2+16-32a+16a^2
\]
即 \(16a^2+128a-159=0\). 因而
\[
a=\frac{-128\pm\sqrt{128^2+4\cdot16\cdot159}}{32}
=\frac{-128\pm8\sqrt{415}}{32}
=\frac{-16\pm\sqrt{415}}4
\]
由于 \(20<\sqrt{415}<21\)，正根 \(a_+=\frac{-16+\sqrt{415}}4\) 满足 \(1<a_+<\frac54\)，于是 \(z=2-2a_+<0\)，且 \(y=\frac{15-16a_+}{2}<0\)，不符合直角边为边长的要求. 负根 \(a_-=\frac{-16-\sqrt{415}}4<-9\) 时，\(x=10-8a_->0\)，\(y=\frac{15-16a_-}{2}>0\)，\(z=2-2a_->0\)，且已经满足上面的勾股等式，因此有效. 故 \(a=\frac{-16-\sqrt{415}}4\).
\end{enumerate}
\end{solution}

\begin{problem}
已知：\(\alpha\)，\(\beta\) 皆为锐角，求证：\(\sin(\alpha+\beta)\leq\sin\alpha+\sin\beta\).
\end{problem}

\begin{solution}
由两角和的正弦公式，
\[
\sin\left(\alpha+\beta\right)=\sin\alpha\cos\beta+\cos\alpha\sin\beta
\]
因为 \(\alpha\)，\(\beta\) 都是锐角，所以 \(\sin\alpha>0\)，\(\sin\beta>0\)，且 \(0<\cos\alpha<1\)，\(0<\cos\beta<1\). 因此
\[
\sin\alpha+\sin\beta-\sin\left(\alpha+\beta\right)
=\sin\alpha\left(1-\cos\beta\right)+\sin\beta\left(1-\cos\alpha\right)>0
\]
这里 \(\sin\alpha\)，\(\sin\beta\)，\(1-\cos\alpha\)，\(1-\cos\beta\) 都为正数，所以
\[
\sin\left(\alpha+\beta\right)<\sin\alpha+\sin\beta
\]
严格不等式必然推出题中的不等式，故命题得证.
\end{solution}

\begin{problem}
已知，正方形 \(ABCD\) 边长为 \(a\)，连接相邻各边中点，得一正方形 \(A_1B_1C_1D_1\)，再连接 \(A_1B_1C_1D_1\) 相邻各边中点又得一个正方形 \(A_2B_2C_2D_2\)，如此相继下去，可得无数多个正方形. 求这些正方形面积的总和.
\end{problem}

\begin{solution}
设第 \(n\) 个正方形的边长为 \(l_n\)，面积为 \(S_n\)，其中原正方形是第一个正方形. 在边长为 \(l_n\) 的正方形中，相邻两边中点的距离为
\[
l_{n+1}=\sqrt{\left(\frac{l_n}{2}\right)^2+\left(\frac{l_n}{2}\right)^2}=\frac{\sqrt2}{2}l_n
\]
所以相邻两个正方形的面积满足 \(S_{n+1}=\frac12S_n\)，且 \(S_1=a^2\). 前 \(N\) 个正方形的面积和为
\[
S_1+S_2+\cdots+S_N=a^2\frac{1-\left(\frac12\right)^N}{1-\frac12}=2a^2\left(1-\left(\frac12\right)^N\right)
\]
令 \(N\) 无限增大，\(\left(\frac12\right)^N\) 趋于 \(0\)，故这些正方形面积的总和为 \(2a^2\).
\end{solution}

\begin{problem}
解不等式：\(2\cos^2x-\left(2+\sqrt3\right)\cos x+\sqrt3<0\).
\end{problem}

\begin{solution}
令 \(t=\cos x\)，则 \(-1\le t\le1\)，并且
\[
2t^2-\left(2+\sqrt3\right)t+\sqrt3=(t-1)(2t-\sqrt3)
\]
因此原不等式等价于 \(\left(t-1\right)\left(2t-\sqrt3\right)<0\). 在 \(-1\le t\le1\) 内，始终有 \(t-1\le0\). 要使乘积为负，必须同时满足 \(t-1<0\) 和 \(2t-\sqrt3>0\)，即
\[
\frac{\sqrt3}{2}<\cos x<1
\]
由 \(\cos x>\frac{\sqrt3}{2}\)，得 \(2k\pi-\frac\pi6<x<2k\pi+\frac\pi6\)，其中 \(k\in\Z\). 条件 \(\cos x<1\) 排除各区间中心的 \(x=2k\pi\). 所以原不等式的解集为
\[
x\in\bigcup_{k\in\Z}\left[\left(2k\pi-\frac\pi6,2k\pi\right)\cup\left(2k\pi,2k\pi+\frac\pi6\right)\right]
\]
\end{solution}
